%%%%%%%%%%%%%%%%%%%%%%%%%%%%%%%%%%%%%%%%%%%%%%%%%%%%%%%%%%%%%%%%%%%%%%
% Lớp 11 | Hàm số lượng giác và phương trình lượng giác | Các công thức lượng giác
% Dạng 1D1?3-4: Áp dụng công thức biến đổi tích <-> tổng
% Số câu: 43  (sinh tự động bởi tach_cau.py; từ nay thầy sửa trực tiếp file này)
%%%%%%%%%%%%%%%%%%%%%%%%%%%%%%%%%%%%%%%%%%%%%%%%%%%%%%%%%%%%%%%%%%%%%%

%%%%%%%%%%%%%%%% TRẮC NGHIỆM 4 PHƯƠNG ÁN %%%%%%%%%%%%%%%%

% nguồn: Ok_Toan11_TracNghiem_Lan2
\begin{ex}%[1D1N3-4]%[ma:3b0fc677]
	Rút gọn biểu thức $A = \dfrac{1 + \sin 4\alpha - \cos 4\alpha}{1 + \sin 4\alpha  + \cos 4\alpha}$.
	\choice 
		{$\sin  2\alpha$}
		{$\cos 2\alpha$} 
		{\True $\tan 2\alpha$} 
		 {$\cot 2\alpha$}
	\loigiai{ 
		Ta có  \\
		$A = \dfrac{(1-\cos 4\alpha)+\sin 4\alpha}{(1+\cos 4\alpha)+\sin 4\alpha} = \dfrac{2\sin^2 2\alpha+2\sin  2\alpha\cos 2\alpha}{2\cos^2 2\alpha+2\sin 2\alpha\cos 2\alpha} = \dfrac{2\sin 2\alpha(\sin 2\alpha+\cos 2\alpha)}{2\cos 2\alpha(\sin 2\alpha+\cos 2\alpha)} = \tan 2\alpha$.
	} 
\end{ex}

% nguồn: Ok_Toan11_TracNghiem_Lan2
\begin{ex}%[1D1N3-4]%[ma:29788ef5]
	Mệnh đề nào  sau đây \textbf{sai}?
	 \choice
		{$\cos a\cos b=\dfrac  {1}{2}\left[\cos (a-b)+\cos (a+b)\right]$}
		{\True $\sin  a\cos b=\dfrac {1}{2}\left[\sin (a-b)-\cos (a+b)\right]$}
		 {$\sin a\sin b=\dfrac {1}{2}\left[\cos (a-b)-\cos (a+b)\right]$}
		{$\sin a\cos b=\dfrac {1}{2}\left[\sin (a-b)+\sin (a+b)\right]$} 
	\loigiai{Ta có  $\sin a\cos b=\dfrac {1}{2}\left[\sin (a+b)+\sin (a-b)\right]$.
	 }
\end{ex}

% nguồn: Ok_Toan11_TracNghiem_Lan2
\begin{ex}%[1D1N3-4]%[ma:bff7b32d]
	Công thức nào sau đây là  \textbf{sai}?
	\choice 
		{$\cos a+\cos b=2\cos  \dfrac {a+b}{2}\cdot \cos \dfrac {a-b}{2}$}
		{$\cos  a-\cos b=-2\sin \dfrac {a+b}{2}\cdot \sin \dfrac {a-b}{2}$}
		{$\sin  a+\sin b=2\sin \dfrac {a+b}{2}\cdot \cos \dfrac {a-b}{2}$}
		{\True $\sin a-\sin b=2\sin \dfrac  {a+b}{2}\cdot \cos \dfrac {a-b}{2}$}
	\loigiai{Ta có  $\sin a-\sin b=2\cos \dfrac {a+b}{2}\cdot \sin \dfrac {a-b}{2}$.
	} 
\end{ex}

% nguồn: Ok_Toan11_TracNghiem_Lan2
\begin{ex}%[1D1N3-4]%[ma:b8505dcf]
	Trong các đẳng thức  sau, đẳng thức nào \textbf{sai}?
	\choice 
		 {$\cos(a-b)=\cos a\cdot \cos b+\sin a\cdot \sin b$}
		{$\cos a\cdot \cos  b=\dfrac {1}{2}\left[\cos(a+b)+\cos(a-b)\right]$}
		{$\sin (a-b)=\sin a\cdot \cos b-\sin  b\cdot \cos a$}
		{\True $\cos a+\cos b=2\cos(a+b)\cdot \cos(a-b)$} 
	\loigiai{Ta có $\cos a+\cos b=2\cos\dfrac {a+b}{2}\cdot \cos\dfrac {a-b}{2}$ 
	} 
\end{ex}

% nguồn: Ok_TracNghiemLop11_Lan1
\begin{ex}%[1D1N3-4]%[ma:86318a06]
	Kết quả  $\cos \left( \dfrac{\pi}{2}-\alpha \right)$ bằng
	 \choice
		 {$\cos \alpha$}
		{\True $\sin  \alpha$}
		{$-\cos  \alpha$}
		 {$-\sin \alpha$}
	 \loigiai{
		 Ta có
		 $\cos \left( \dfrac{\pi}{2}-\alpha \right)=\sin \alpha$
	} 
\end{ex}

% nguồn: Ok_Toan11_TracNghiem_Lan2
\begin{ex}%[1D1H3-4]%[ma:c43c7801]
	Rút gọn biểu thức $A=\dfrac {\sin 3x+\cos 2x-\sin x}{\cos x+\sin 2x-\cos 3x}$ ($\sin 2x\ne 0$; $2\sin x+1\ne  0$) ta được
	 \choice
		{$A=\cot  6x$}
		{$A=\cot 3x$} 
		{\True  $A=\cot 2x$}
		 {$A=\tan x+\tan 2x+\tan 3x$}
	\loigiai{ 
		\allowdisplaybreaks 
		\begin{align*}
			A&=\dfrac {\sin 3x+\cos 2x-\sin x}{\cos x+\sin 2x-\cos  3x}\\
			&=\dfrac {2\cos 2x\sin x+\cos 2x}{2\sin  2x\sin x+\sin 2x}\\
			&=\dfrac {\cos 2x(1+2\sin x)}{\sin 2x(1+2\sin x)}\\ 
			&=\cot 2x. 
		\end{align*}
	} 
\end{ex}

% nguồn: Ok_Toan11_TracNghiem_Lan2
\begin{ex}%[1D1H3-4]%[ma:0f091a04]
	Tích số $\cos 10^{\circ}\cdot\cos 30^{\circ}\cdot \cos 50^{\circ}\cdot \cos 70^{\circ}$ bằng 
	\choice 
		 {$\dfrac{1}{16}$}
		 {$\dfrac{1}{8}$}
		{\True  $\dfrac{3}{16}$}
		 {$\dfrac{1}{4}$}
	\loigiai{ 
		 Ta có
		\begin{eqnarray*}
			&&\cos 10^{\circ}\cdot \cos 30^{\circ}\cdot \cos 50^{\circ}\cdot \cos 70^{\circ}\\ 
			&=&\cos 10^{\circ}\cdot\cos  30^{\circ}\cdot\dfrac{1}{2} \left(\cos 120^\circ+\cos 20^\circ \right)\\
			&=&\dfrac{\sqrt{3}}{4} \left(-\dfrac{\cos 10^{\circ}}{2}+\dfrac{\cos  30^{\circ}+\cos 10^{\circ}}{2} \right)\\
			 &=&\dfrac{\sqrt{3}}{4}\cdot\dfrac{\sqrt{3}}{4}=\dfrac{3}{16}.
		\end{eqnarray*}
	} 
\end{ex}

% nguồn: Ok_Toan11_TracNghiem_Lan2
\begin{ex}%[1D1H3-4]%[ma:4048431e]
	Rút gọn $M  = \cos \left(x + \dfrac{\pi}{4}\right) - \cos \left(x - \dfrac{\pi}{4}\right)$.
	 \choice
		{$M = \sqrt{2}  \sin x$}
		{\True  $M = -\sqrt{2} \sin x$}
		 {$M = \sqrt{2} \cos x$}
		{$M = -\sqrt{2} \cos  x$}
	 \loigiai{
		Áp dụng công thức $\cos  a-\cos b=-2\sin \dfrac{a+b}{2}\cdot\sin \dfrac{a-b}{2}$, ta được\\
		$M=\cos \left( x+\dfrac{\pi }{4} \right)-\cos \left( x-\dfrac{\pi }{4} \right)=-2\sin  \left( \dfrac{x+\dfrac{\pi }{4}+x-\dfrac{\pi }{4}}{2} \right)\cdot\sin \left( \dfrac{x+\dfrac{\pi }{4}-x+\dfrac{\pi }{4}}{2} \right)$\\
		$=-2\sin  x\cdot\sin \dfrac{\pi }{4}=-\sqrt{2}\sin x$.
	 }
\end{ex}

% nguồn: Ok_Toan11_TracNghiem_Lan2
\begin{ex}%[1D1H3-4]%[ma:5bd235a4]
	Rút gọn biểu thức  $A=\dfrac{\sin x+\sin 2x+\sin 3x}{\cos x+\cos 2x+\cos 3x}$.
	\choice 
		{$A=\tan 6x$} 
		{$A=\tan  3x$}
		 {\True $A=\tan 2x$}
		{$A=\tan x+\tan  2x+\tan 3x$}
	 \loigiai{
		 Ta có
		$A=\dfrac{\sin x+\sin 2x+\sin 3x}{\cos x+\cos 2x+\cos 3x}=\dfrac{2\sin 2x\cdot\cos x+\sin 2x}{2\cos 2x\cdot\cos x+\cos 2x}=\dfrac{\sin 2x\left(2\cos x+1\right)}{\cos 2x\left(2\cos  x+1\right)}=\tan 2x$.
	 }
\end{ex}

% nguồn: Ok_Toan11_TracNghiem_Lan2
\begin{ex}%[1D1H3-4]%[ma:8c7062f1]
	Biến đổi biểu thức $\sin \alpha-1$ thành tích. 
	 \choice
		{$\sin \alpha -1=2\sin \left( \alpha -\dfrac {\pi }{2} \right)\cos \left( \alpha  +\dfrac {\pi }{2} \right)$}
		{\True $\sin \alpha -1=2\sin \left( \dfrac {\alpha }{2}-\dfrac {\pi }{4} \right)\cos \left( \dfrac {\alpha }{2}+\dfrac  {\pi }{4} \right)$}
		{$\sin \alpha  -1=2\sin \left(\alpha +\dfrac {\pi }{2} \right)\cos \left( \alpha -\dfrac {\pi }{2} \right)$}
		{$\sin \alpha -1=2\sin  \left( \dfrac {\alpha }{2}+\dfrac {\pi }{4} \right)\cos \left( \dfrac {\alpha }{2}-\dfrac {\pi }{4} \right)$}
	 \loigiai{\allowdisplaybreaks
		\begin{align*}
			\sin \alpha -1&=\sin \alpha -\sin \dfrac  {\pi }{2}\\
			&=2\cos \dfrac {\alpha +\dfrac {\pi }{2}}{2}\sin \dfrac {\alpha -\dfrac {\pi  }{2}}{2}\\
			&=2\cos \left( \dfrac {\alpha }{2}+\dfrac {\pi }{4} \right)\sin \left( \dfrac {\alpha  }{2}-\dfrac {\pi }{4} \right).
		\end{align*}
	} 
\end{ex}

% nguồn: Ok_Toan11_TracNghiem_Lan2
\begin{ex}%[1D1H3-4]%[ma:35247150]
	Tích số ${\cos  \dfrac{\pi }{7}\cdot \cos \dfrac{4\pi }{7}\cdot \cos \dfrac{5\pi }{7}}$ bằng
	\choice 
		{\True $\dfrac  {1}{8}$}
		 {$-\dfrac {1}{8}$}
		{$\dfrac  {1}{4}$}
		{$-\dfrac {1}{4}$} 
	 \loigiai{
		\allowdisplaybreaks 
		\begin{align*}
			\cos \dfrac {\pi }{7}\cdot  \cos \dfrac {4\pi }{7}\cdot \cos \dfrac {5\pi }{7}&=\dfrac {\sin \dfrac {2\pi }{7}\cdot \cos \dfrac {4\pi }{7}\cdot \cos \dfrac {5\pi }{7}}{2\sin \dfrac {\pi }{7}}\\
			&=-\dfrac {\sin \dfrac {2\pi }{7}\cdot \cos \dfrac {2\pi }{7}\cdot \cos \dfrac {4\pi  }{7}}{2\sin \dfrac {\pi }{7}}\\
			&=-\dfrac {\sin \dfrac {4\pi }{7}\cdot  \cos \dfrac {4\pi }{7}}{4\sin \dfrac {\pi }{7}}\\
			&=-\dfrac {\sin \dfrac  {8\pi }{7}}{8\sin \dfrac {\pi }{7}}\\
			&=\dfrac {1}{8}. 
		\end{align*}
	 }
\end{ex}

% nguồn: Ok_Toan11_TracNghiem_Lan2
\begin{ex}%[1D1H3-4]%[ma:c9677f10]
	 Biểu thức $A=\dfrac {1}{2\sin 10^\circ}-2\sin 70^\circ$ có giá trị đúng bằng
	\choice 
		{\True  $1$}
		{$1$} 
		{$2$} 
		 {$3$}
	 \loigiai{
		\allowdisplaybreaks 
		\begin{align*}
			A&=\dfrac {1}{2\sin  10^\circ}-2\sin 70^\circ\\
			&=\dfrac {1-4\sin 10^\circ \cdot \sin 70^\circ}{2\sin  10^\circ}\\
			 &=\dfrac {2\sin 80^\circ}{2\sin 10^\circ}\\
			&=\dfrac  {2\sin 10^\circ}{2\sin 10^\circ}\\
			&=1. 
		\end{align*}
	} 
\end{ex}

% nguồn: Ok_Toan11_TracNghiem_Lan2
\begin{ex}%[1D1H3-4]%[ma:e4d706f0]
	Gọi $M = \cos x  + \cos 2x + \cos 3x$. Rút gọn biểu thức $M$, ta được
	 \choice
		{$M = 2 \cos 2x (\cos x +  1)$}
		{$M = 4 \cos 2x \left(\dfrac{1}{2} + \cos  x\right)$}
		{$M =  \cos 2x (2 \cos x - 1)$}
		{\True $M = \cos 2x (2  \cos x + 1)$}
	 \loigiai{
		Ta có $M = \cos x + \cos 2x + \cos 3x = (\cos x + \cos 3x) +  \cos 2x= 2\cos 2x\cdot\cos x + \cos 2x = \cos 2x(2\cos x+1)$.
	} 
\end{ex}

% nguồn: Ok_Toan11_TracNghiem_Lan2
\begin{ex}%[1D1H3-4]%[ma:66f3d0da]
	Số  khẳng định đúng trong các khẳng định sau:\\
	 $(I)$ $\cos a\cos b=\dfrac {1}{2}\left[\cos (a-b)+\cos (a+b)\right]$.\\
	$(II)$ $\sin a\sin b=\dfrac {1}{2}\left[\cos (a-b)-\cos (a+b)\right]$.\\ 
	$(III)$ $\cos a+\cos b=2\cos \dfrac {a+b}{2}\cos \dfrac  {a-b}{2}$.\\
	$(IV)$ $\sin a-\sin b=2\cos \dfrac {a+b}{2}\cos  \dfrac {a-b}{2}$.
	 \choice
		 {$0$}
		{$1$} 
		{$2$} 
		{\True  $3$}
	\loigiai{Khẳng định $(IV)$ sai nên có $3$ khẳng  định đúng.
	} 
\end{ex}

% nguồn: Ok_Toan11_TracNghiem_Lan2
\begin{ex}%[1D1H3-4]%[ma:2d0b76e3]
	Cho hai góc nhọn $a$ và $b$. Biết $\cos a=\dfrac{1}{3}$, $\cos b=\dfrac{1}{4}$. Giá trị  $\cos\left(a+b\right)\cdot\cos\left(a-b\right)$ bằng
	 \choice
		{$-\dfrac{113}{144}$} 
		 {$-\dfrac{115}{144}$}
		{$-\dfrac{117}{144}$} 
		 {\True $-\dfrac{119}{144}$}
	 \loigiai{
		Ta có \begin{align*}\cos\left(a+b\right)\cdot\cos\left(a-b\right)&=\dfrac{1}{2}\left(\cos 2a+\cos 2b\right)\\&=\cos^2a+\cos ^2b-1\\
		&=\left(\dfrac{1}{3}\right)^2+\left(\dfrac{1}{4}\right)^2-1\\&=-\dfrac{119}{144}. 
	\end{align*}
 }
\end{ex}

% nguồn: Ok_Toan11_TracNghiem_Lan2
\begin{ex}%[1D1H3-4]%[ma:6126643a]
	Trong các khẳng định sau, khẳng định  nào \textbf{sai}?
	\choice 
		{$\cos 3x+\cos x=2\cos 2x\cdot  \cos x$}
		{\True $\cos 3x-\cos  x=2\sin 2x\cdot \sin x$}
		{$\sin 3x-\sin  x=2\cos 2x\cdot \sin x$}
		{$\sin 3x+\sin x=2\sin  2x\cdot \cos x$}
	\loigiai{$\cos 3x-\cos x=-2\sin 2x\cdot \sin  x$.
	 }
\end{ex}

% nguồn: Ok_Toan11_TracNghiem_Lan2
\begin{ex}%[1D1H3-4]%[ma:65ec5ace]
	Giá trị đúng của $\cos \dfrac {2\pi }{7}+\cos \dfrac {4\pi  }{7}+\cos \dfrac {6\pi }{7}$ bằng
	\choice 
		 {$\dfrac {1}{2}$}
		{\True $-\dfrac  {1}{2}$}
		{$\dfrac {1}{4}$} 
		 {$-\dfrac {1}{4}$}
	\loigiai{Ta có 
		\allowdisplaybreaks 
		\begin{align*}
			\cos \dfrac {2\pi }{7}+\cos \dfrac {4\pi }{7}+\cos \dfrac  {6\pi }{7}&=\dfrac {\sin \dfrac {\pi }{7}\left( \cos \dfrac {2\pi }{7}+\cos \dfrac {4\pi }{7}+\cos \dfrac {6\pi }{7} \right)}{\sin \dfrac {\pi }{7}}\\
			&=\dfrac {\sin \dfrac {3\pi }{7}+\sin \left( -\dfrac {\pi }{7} \right)+\sin \dfrac {5\pi  }{7}+\sin \left( -\dfrac {3\pi }{7} \right)+\sin \pi +\sin \left( -\dfrac {5\pi }{7} \right)}{2\sin \dfrac {\pi }{7}}\\
			&=\dfrac  {\sin \left( -\dfrac {\pi }{7} \right)}{2\sin \dfrac {\pi }{7}}\\
			 &=-\dfrac {1}{2}.
		\end{align*}
	 }
\end{ex}

% nguồn: Ok_Toan11_TracNghiem_Lan2
\begin{ex}%[1D1H3-4]%[ma:137ca5f3]
	Rút gọn biểu thức $A = \dfrac{1 + \cos x + \cos 2x + \cos 3x}{2 \cos^2 x + \cos x -  1}$.
	 \choice
		 {$\cos x$}
		{$2 \cos  x - 1$}
		{\True $2 \cos  x$}
		{$\cos x - 1$} 
	\loigiai{ 
		Ta có $A = \dfrac{(1+\cos 2x)+(\cos x+\cos 3x)}{(2\cos^2 x-1)+\cos x} = \dfrac{2\cos^2 x+2\cos 2x\cos x}{\cos x+\cos 2x}= \dfrac{2\cos x(\cos x+\cos  2x)}{\cos x+\cos 2x} = 2\cos x$.
	} 
\end{ex}

% nguồn: Ok_Toan11_TracNghiem_Lan2
\begin{ex}%[1D1H3-4]%[ma:2256b741]
	Với $\alpha$ là số thực bất  kỳ, mệnh đề nào sau đây là mệnh đề đúng?
	 \choice
		{$\cos 2\alpha+\cos 4a=2\cos 2\alpha \cdot  \cos 6\alpha$}
		{$\sin 2\alpha+\sin 4a=2\sin  \alpha \cdot \cos 3\alpha$}
		{$\cos 2\alpha-\cos 4a=-2\sin  3\alpha \cdot \sin \alpha$}
		 {\True $\sin 2\alpha-\sin 4a=-2\cos 3\alpha \cdot \sin \alpha$}
	\loigiai{Ta  có\\
		$\cos 2\alpha+\cos  4a=2\cos \dfrac {2\alpha+4\alpha}{2}\cos \dfrac {2\alpha-4\alpha }{2}=2\cos 3\alpha \cos \alpha$.\\
		$\sin 2\alpha +\sin 4a=2\sin \dfrac {2\alpha +4\alpha }{2}\cdot \cos \dfrac  {2\alpha -4\alpha }{2}=2\sin 3\alpha \cdot \cos \alpha$.\\
		$\cos 2\alpha -\cos 4a=-2\sin \dfrac {2\alpha +4\alpha }{2}\cdot \sin \dfrac {2\alpha -4\alpha }{2}=2\sin 3\alpha  \cdot \sin \alpha$.\\
		$\sin 2\alpha -\sin 4a=2\cos \dfrac {2\alpha +4\alpha }{2}\cdot \sin \dfrac  {2\alpha -4\alpha }{2}=-2\cos 3\alpha \cdot \sin \alpha$.
	 }
\end{ex}

% nguồn: Ok_Toan11_TracNghiem_Lan2
\begin{ex}%[1D1H3-4]%[ma:c1355ba1]
	Tính giá trị biểu thức $P=\sin 30^\circ \cdot \cos 60^\circ+\sin  60^\circ \cdot \cos 30^\circ$.
	 \choice
		{\True  $P=1$}
		 {$P=0$}
		{$P=\sqrt  {3}$}
		{$P=-\sqrt {3}$} 
	\loigiai{Ta có $P=\sin  \left(30^\circ+60^\circ\right)=\sin 90^\circ=1$.
	} 
\end{ex}

% nguồn: Ok_Toan11_TracNghiem_Lan2
\begin{ex}%[1D1H3-4]%[ma:2aabb898]
	Rút gọn biểu  thức $P=\sin \left(a+\dfrac {\pi }{4} \right)\sin \left( a-\dfrac {\pi }{4} \right)$.
	 \choice
		{$-\dfrac {3}{2}\cos  2a$}
		{$\dfrac {1}{2}\cos 2a$} 
		{$-\dfrac  {2}{3}\cos 2a$}
		{\True  $-\dfrac {1}{2}\cos 2a$}
	\loigiai{Ta có $\sin \left(a+\dfrac {\pi }{4}  \right)\sin \left(a-\dfrac {\pi }{4} \right)=\dfrac {1}{2}\left[\cos \dfrac {\pi }{2}-\cos 2a \right]=-\dfrac {1}{2}\cos 2a$.
	} 
\end{ex}

% nguồn: Ok_Toan11_TracNghiem_Lan2
\begin{ex}%[1D1H3-4]%[ma:da75e022]
	Rút gọn biểu thức $\cos \left(x + \dfrac{\pi}{4}\right)  - \cos \left(x - \dfrac{\pi}{4}\right)$, ta được
	\choice 
		{$\sqrt{2}  \sin x$}
		 {\True $-\sqrt{2} \sin x$}
		{$\sqrt{2}  \cos x$}
		{$-\sqrt{2} \cos  x$}
	\loigiai{ 
		Ta có $$\cos \left(x + \dfrac{\pi}{4}\right) - \cos \left(x - \dfrac{\pi}{4}\right) =  -2\sin \left(\dfrac{x + \dfrac{\pi}{4} + x - \dfrac{\pi}{4}}{2}\right) \sin \left(\dfrac{x + \dfrac{\pi}{4} - x +\dfrac{\pi}{4}}{2}\right)= -2\sin x \sin \dfrac{\pi}{4} = -\sqrt{2} \sin x.$$
	 }
\end{ex}

% nguồn: Ok_Toan11_TracNghiem_Lan2
\begin{ex}%[1D1H3-4]%[ma:df3528cc]
	Giá trị đúng của $\tan \dfrac  {\pi }{24}+\tan \dfrac {7\pi }{24}$ bằng
	 \choice
		 {\True $2\left( \sqrt {6}-\sqrt {3} \right)$}
		{$2\left(\sqrt {6}+\sqrt {3} \right)$} 
		{$2\left(\sqrt {3}-\sqrt  {2} \right)$}
		 {$2\left(\sqrt {3}+\sqrt {2} \right)$}
	\loigiai{ 
		 \allowdisplaybreaks
		\begin{align*}
			\tan \dfrac  {\pi }{24}+\tan \dfrac {7\pi }{24}&=\dfrac {\sin \dfrac {\pi }{3}}{\cos \dfrac {\pi }{24}\cdot \cos \dfrac {7\pi }{24}}\\
			&=\dfrac {\sqrt {3}}{\cos \dfrac {\pi }{3}+\cos \dfrac {\pi }{4}}\\ 
			 &=2\left(\sqrt {6}-\sqrt {3} \right).
		\end{align*}
	 }
\end{ex}

% nguồn: Ok_Toan11_TracNghiem_Lan2
\begin{ex}%[1D1H3-4]%[ma:ebc4f2ab]
	Rút gọn biểu thức $P=\dfrac {\cos a+2\cos 3a+\cos 5a}{\sin a+2\sin  3a+\sin 5a}$.
	 \choice
		{$P=\tan a$} 
		{$P=\cot  a$}
		{\True  $P=\cot 3a$}
		 {$P=\tan 3a$}
	 \loigiai{
		 \allowdisplaybreaks
		\begin{align*}
			P&=\dfrac {\cos a+2\cos 3a+\cos  5a}{\sin a+2\sin 3a+\sin 5a}\\
			&=\dfrac {2\cos 3a\cos a+2\cos  3a}{2\sin 3a\cos a+2\sin 3a}\\
			&=\dfrac  {2\cos 3a( \cos a+1)}{2\sin 3a( \cos a+1)}\\
			&=\dfrac {\cos 3a}{\sin 3a}\\ 
			&=\cot  3a.
		\end{align*}
	} 
\end{ex}

% nguồn: Ok_Toan11_TracNghiem_Lan2
\begin{ex}%[1D1H3-4]%[ma:d390d85f]
	Tích số $\cos 10^\circ \cdot \cos 30^\circ \cdot \cos 50^\circ \cdot \cos 70^\circ$  bằng
	\choice 
		{$\dfrac  {1}{16}$}
		 {$\dfrac {1}{8}$}
		{\True  $\dfrac {3}{16}$}
		{$\dfrac {1}{4}$} 
	 \loigiai{\allowdisplaybreaks
		\begin{align*}
			\cos  10^\circ \cdot \cos 30^\circ \cdot \cos 50^\circ \cdot \cos 70^\circ &=\cos 10^\circ \cdot \cos 30^\circ \cdot \dfrac {1}{2}\left( \cos 120^\circ+\cos 20^\circ \right)\\
			&=\dfrac {\sqrt {3}}{4}\left( -\dfrac {\cos 10^\circ  }{2}+\dfrac {\cos 30^\circ +\cos 10^\circ }{2} \right)\\
			 &=\dfrac {\sqrt {3}}{4}\cdot \dfrac {\sqrt{3}}{4}\\
			&=\dfrac {3}{16}. 
		\end{align*}
	 }
\end{ex}

% nguồn: Ok_Toan11_TracNghiem_Lan2
\begin{ex}%[1D1H3-4]%[ma:533758cb]
	Rút gọn biểu  thức $M = \dfrac{\sin 3x - \sin x}{2 \cos^2 x - 1}$.
	\choice 
		{$\tan  2x$}
		 {\True $2 \sin x$}
		{$2 \tan  x$}
		 {$\sin 2x$}
	 \loigiai{
		Ta có $\dfrac{\sin 3x - \sin x}{2\cos^2 x - 1} = \dfrac{2\cos 2x  \sin x}{\cos 2x} = 2\sin x$.
	} 
\end{ex}

% nguồn: Ok_TracNghiemLop11_Lan1
\begin{ex}%[1D1H3-4]%[ma:e001cecb]
	Rút gọn các biểu thức  $A=\sin\left(\dfrac{\pi}{2}+\alpha\right)+\sin\left(\dfrac{\pi}{2}-\alpha\right)$.
	 \choice
		 {$A=2$}
		{\True  $A=2\cos(\alpha)$}
		 {$A=0$}
		 {$A=2\sin(\alpha)$}
	 \loigiai{
		\begin{eqnarray*}
			A &=&  \sin \left( {\dfrac{\pi }{2} + \alpha } \right) + \sin \left( {\dfrac{\pi }{2} - \alpha } \right)\\
			&=& 2\sin \left[{\dfrac{{\left( {\dfrac{\pi  }{2} + \alpha } \right) + \left( {\dfrac{\pi }{2} - \alpha } \right)}}{2}} \right]\cos \left[{\dfrac{{\left( {\dfrac{\pi }{2} + \alpha } \right) - \left( {\dfrac{\pi }{2} - \alpha } \right)}}{2}} \right]\\
			&=& 2\sin \left({\dfrac{\pi }{2}} \right)\cos \left(  {\dfrac{{2\alpha }}{2}} \right) \\
			&=& 2\cos \left( {\alpha  } \right).
		\end{eqnarray*}
	} 
\end{ex}

% nguồn: Ok_Toan11_TracNghiem_Lan2
\begin{ex}%[1D1V3-4]%[ma:93bb48a6]
	Biến đổi  biểu thức $\sin a+1$ thành tích.
	\choice 
		{$\sin a+1=2\sin\left(\dfrac{a}{2}+\dfrac{\pi}{4}\right)\cos\left(\dfrac{a}{2}-\dfrac{\pi}{4}\right)$} 
		{$\sin a+1=2\cos\left(\dfrac{a}{2}+\dfrac{\pi}{4}\right)\sin\left(\dfrac{a}{2}-\dfrac{\pi}{4}\right)$} 
		{$\sin  a+1=2\sin\left(a+\dfrac{\pi}{2}\right)\cos\left(a-\dfrac{\pi}{2}\right)$}
		 {\True $\sin a+1=2\cos\left(a+\dfrac{\pi}{2}\right)\sin\left(a-\dfrac{\pi}{2}\right)$}
	\loigiai{ 
		Ta có \begin{align*}\sin a+1&=2\sin\dfrac{a}{2}\cos\dfrac{a}{2}+\sin ^2\dfrac{a}{2}+\cos ^2\dfrac{a}{2}\\
		 &=\left(\sin\dfrac{a}{2}+\cos\dfrac{a}{2}\right)^2\\&=2\sin ^2\left(\dfrac{a}{2}+\dfrac{\pi}{4}\right)\\
		 &=2\sin\left(\dfrac{a}{2}+\dfrac{\pi}{4}\right)\cos\left(\dfrac{\pi}{4}-\dfrac{a}{2}\right)\\
		&=2\sin\left(\dfrac{a}{2}+\dfrac{\pi}{4}\right)\cos\left(\dfrac{a}{2}-\dfrac{\pi}{4}\right). 
	\end{align*}
 }
\end{ex}


%%%%%%%%%%%%%%%% ĐÚNG/SAI %%%%%%%%%%%%%%%%

% nguồn: Ok_dungsai_lop11_Lan2
\begin{ex}%[1D1N3-4]%[ma:8699e0fc]
	Xét tính đúng sai của các mệnh đề sau 
	\choiceTF 
		{\True $\sin a+\sin  b=2\cdot\sin\dfrac{a+b}{2}\cdot\sin\dfrac{a-b}{2}$}
		{\True $\cos a+\cos  b=2\cdot\cos\dfrac{a+b}{2}\cos\cdot\dfrac{a-b}{2}$}
		{$\cos a-\cos  b=2\cdot\sin\dfrac{a+b}{2}\cdot\sin\dfrac{a-b}{2}$}
		{\True  $\sin a-\sin b=2\cdot\cos\dfrac{a+b}{2}\cdot\sin\dfrac{a-b}{2}$}
	 \loigiai{
	\begin{itemchoice}
		\itemch  \textbf{Đúng.} \\
		 Vì $\sin a+\sin b=2\cdot\sin\dfrac{a+b}{2}\cdot\sin\dfrac{a-b}{2}$.
		\itemch \textbf{Đúng.}  \\
		Vì  $\cos a+\cos b=2\cdot\cos\dfrac{a+b}{2}\cdot\cos\dfrac{a-b}{2}$.
		\itemch \textbf{Sai.}  \\
		 Vì $\cos a-\cos b=-2\cdot\sin\dfrac{a+b}{2}\cdot\sin\dfrac{a-b}{2}$
		\itemch  \textbf{Đúng.} \\
		Vì $\sin a-\sin  b=2\cdot\cos\dfrac{a+b}{2}\cdot\sin\dfrac{a-b}{2}$.
	\end{itemchoice}
	 }
\end{ex}

% nguồn: Ok_dungsai_lop11_Lan1
\begin{ex}%[1D1H3-4]%[ma:044705b5]
 Mỗi kết quả  dưới đây đúng hay sai?
  \choiceTF
  {\True  $\sin\left(\dfrac{\pi}{3}+\alpha\right)\sin\left(\dfrac{\pi}{3}-\alpha\right)=\dfrac{\cos2\alpha}{2}+\dfrac{1}{4}$}
  {\True $\sin\left(\dfrac{\pi}{3}+\alpha\right)\cos\left(\dfrac{\pi}{3}-\alpha\right)=\dfrac{\sin2\alpha}{2}+\dfrac{\sqrt{3}}{4}$} 
   {$\cos\left(\dfrac{\pi}{3}+\alpha\right)\cos\left(\dfrac{\pi}{3}-\alpha\right)=-\dfrac{\cos2\alpha}{2}-\dfrac{1}{4}$}
  {$\cos\left(\dfrac{\pi}{3}+\alpha\right)\sin\left(\dfrac{\pi}{3}-\alpha\right)=\dfrac{\sin2\alpha}{2}-\dfrac{\sqrt{3}}{4}$} 
  \loigiai{
  \begin{itemchoice}
   \itemch $\sin\left(\dfrac{\pi}{3}+\alpha\right)\sin\left(\dfrac{\pi}{3}-\alpha\right)=\dfrac{1}{2}\left(\cos2\alpha-\cos\dfrac{2\pi}{3}\right)=\dfrac{\cos2\alpha}{2}+\dfrac{1}{4}$. 
   \itemch $\sin\left(\dfrac{\pi}{3}+\alpha\right)\cos\left(\dfrac{\pi}{3}-\alpha\right)=\dfrac{1}{2}\left(\sin2\alpha+\sin\dfrac{2\pi}{3}\right)=\dfrac{\sin2\alpha}{2}+\dfrac{\sqrt{3}}{4}$. 
    \itemch $\cos\left(\dfrac{\pi}{3}+\alpha\right)\cos\left(\dfrac{\pi}{3}-\alpha\right)=\dfrac{1}{2}\left(\cos2\alpha+\cos\dfrac{2\pi}{3}\right)=\dfrac{\cos2\alpha}{2}-\dfrac{1}{4}$.
    \itemch $\cos\left(\dfrac{\pi}{3}+\alpha\right)\sin\left(\dfrac{\pi}{3}-\alpha\right)=\dfrac{1}{2}\left(\sin\dfrac{2\pi}{3}-\sin2\alpha\right)=\dfrac{\sqrt{3}}{4}-\dfrac{\sin2\alpha}{2}$.
  \end{itemchoice}
  }
\end{ex}

% nguồn: Ok_dungsai_lop11_Lan1
\begin{ex}%[1D1H3-4]%[ma:84b51e41]
 Mỗi  kết quả dưới đây đúng hay sai?
  \choiceTF
   {$\cos\left(\dfrac{\pi}{3}+\alpha\right)\cos\left(\dfrac{\pi}{3}-\alpha\right)=-\dfrac{\cos2\alpha}{2}-\dfrac{1}{4}$}
  {\True  $\sin\left(\dfrac{\pi}{3}+\alpha\right)\sin\left(\dfrac{\pi}{3}-\alpha\right)=\dfrac{\cos2\alpha}{2}+\dfrac{1}{4}$}
  {\True  $\sin\left(\dfrac{\pi}{3}+\alpha\right)\cos\left(\dfrac{\pi}{3}-\alpha\right)=\dfrac{\sin2\alpha}{2}+\dfrac{\sqrt{3}}{4}$}
  {$\cos\left(\dfrac{\pi}{3}+\alpha\right)\sin\left(\dfrac{\pi}{3}-\alpha\right)=\dfrac{\sin2\alpha}{2}-\dfrac{\sqrt{3}}{4}$} 
  \loigiai{
  \begin{itemchoice}
   \itemch $\cos\left(\dfrac{\pi}{3}+\alpha\right)\cos\left(\dfrac{\pi}{3}-\alpha\right)=\dfrac{1}{2}\left(\cos2\alpha+\cos\dfrac{2\pi}{3}\right)=\dfrac{\cos2\alpha}{2}-\dfrac{1}{4}$. 
   \itemch  $\sin\left(\dfrac{\pi}{3}+\alpha\right)\sin\left(\dfrac{\pi}{3}-\alpha\right)=\dfrac{1}{2}\left(\cos2\alpha-\cos\dfrac{2\pi}{3}\right)=\dfrac{\cos2\alpha}{2}+\dfrac{1}{4}$.
   \itemch $\sin\left(\dfrac{\pi}{3}+\alpha\right)\cos\left(\dfrac{\pi}{3}-\alpha\right)=\dfrac{1}{2}\left(\sin2\alpha+\sin\dfrac{2\pi}{3}\right)=\dfrac{\sin2\alpha}{2}+\dfrac{\sqrt{3}}{4}$. 
    \itemch $\cos\left(\dfrac{\pi}{3}+\alpha\right)\sin\left(\dfrac{\pi}{3}-\alpha\right)=\dfrac{1}{2}\left(\sin\dfrac{2\pi}{3}-\sin2\alpha\right)=\dfrac{\sqrt{3}}{4}-\dfrac{\sin2\alpha}{2}$.
  \end{itemchoice}
  }
\end{ex}

% nguồn: Ok_dungsai_lop11_Lan1
\begin{ex}%[1D1H3-4]%[ma:a454edfc]
 Mỗi kết quả dưới đây  đúng hay sai?
  \choiceTF
   {$\sin3x+\sin x=\sin4x$}
  {$\sin3x-\sin x=\cos2x\sin x$} 
   {\True $\cos3x+\cos x=2\cos2x\cos x$}
  {$\cos3x-\cos x=2\sin2x\sin x$} 
  \loigiai{
  \begin{itemchoice}
   \itemch  $\sin3x+\sin x=2\sin\dfrac{3x+x}{2}\cdot\cos\dfrac{3x-x}{2}=2\sin2x\cos x$.
   \itemch $\sin3x-\sin x=2\cos\dfrac{3x+x}{2}\cdot\sin\dfrac{3x-x}{2}=2\cos2x\sin  x$.
   \itemch $\cos3x+\cos x=2\cos\dfrac{3x+x}{2}\cdot\cos\dfrac{3x-x}{2}=2\cos2x\cos x$. 
    \itemch $\cos3x-\cos x=-2\sin\dfrac{3x+x}{2}\cdot\sin\dfrac{3x-x}{2}=-2\sin2x\sin x$.
  \end{itemchoice}
  }
\end{ex}

% nguồn: Ok_dungsai_lop11_Lan2
\begin{ex}%[1D1H3-4]%[ma:bf0857cf]
	Biến đổi được các  biểu thức sau về dạng tích số. Khi đó
	 \choiceTF
		{$\cos 3x+\cos x=2\cos  2x\cdot\cos 3x$}
		{$\sin  3x+\sin 2x=2\sin 2x\cos\dfrac{x}{2}$}
		{\True $\cos  4x-\cos x=-2\sin\dfrac{5x}{2}\sin\dfrac{3x}{2}$}
		{\True $\sin 5x-\sin x=2\cos 3x\sin  2x$}
	\loigiai{ 
		\begin{itemchoice}
			 \itemch \textbf{Sai.} \\
			 $\cos 3x+\cos x=2\cos 2x\cdot\cos x$.
			\itemch  \textbf{Sai.} \\
			$\sin 3x+\sin 2x=2\sin\dfrac{5x}{2}\cos\dfrac{x}{2}$. 
			\itemch  \textbf{Đúng.} \\
			 $\cos 4x-\cos x=-2\sin\dfrac{5x}{2}\sin\dfrac{3x}{2}$.
			\itemch \textbf{Đúng.} \\ 
			$\sin 5x-\sin x=2\cos 3x\sin 2x$. 
		\end{itemchoice}
	 }
\end{ex}

% nguồn: Ok_dungsai_lop11_Lan2
\begin{ex}%[1D1H3-4]%[ma:abb4ad62]
	Cho $\sin{\alpha}=\dfrac{2}{3}, \dfrac{\pi}{2}<\alpha<\pi$. Khi đó 
	\choiceTF 
		 {\True $\cos{\alpha}=-\dfrac{\sqrt{5}}{3}$}
		 {\True $\tan{\alpha}=-\dfrac{2\sqrt{5}}{5}$}
		 {$\cos{\left(\dfrac{\pi}{3}+\alpha\right)}=\dfrac{\sqrt{5}-2\sqrt{3}}{6}$}
		 {$\cos{\left(\dfrac{\pi}{4}-\alpha\right)}=\dfrac{\sqrt{10}-2\sqrt{2}}{6}$}
	\loigiai{ 
	Ta có $\sin^2{\alpha}+\cos^2{\alpha}=1$. Mà $\sin{\alpha}=\dfrac{2}{3}$ nên  suy ra \[\cos^2{\alpha}=1-\sin^2{\alpha}=1-\left(\dfrac{2}{3}\right)^2=\dfrac{5}{9}\Rightarrow \cos{\alpha}=\pm\dfrac{\sqrt{5}}{3}.\]
	Mà $\dfrac{\pi}{2}<\alpha<\pi$  nên suy ra $\cos{\alpha}=-\dfrac{\sqrt{5}}{3}$.\\
	Khi đó  $\tan{\alpha}=\dfrac{\sin{\alpha}}{\cos{\alpha}}=-\dfrac{2\sqrt{5}}{5}$.\\
	Ta  có
	\begin{itemize}
		\item  $\cos{\left(\dfrac{\pi}{3}+\alpha\right)}=\cos{\dfrac{\pi}{3}}\cos{\alpha}-\sin{\dfrac{\pi}{3}}\sin{\alpha}=\dfrac{1}{2}\cdot\left(-\dfrac{\sqrt{5}}{3}\right)-\dfrac{\sqrt{3}}{2}\cdot\dfrac{2}{3}=\dfrac{-\sqrt{5}-2\sqrt{3}}{6}$.
		\item  $\cos{\left(\dfrac{\pi}{4}-\alpha\right)}=\cos{\dfrac{\pi}{4}}\cos{\alpha}\sin{\dfrac{\pi}{4}}\sin{\alpha}=\dfrac{\sqrt{2}}{2}\cdot\left(-\dfrac{\sqrt{5}}{3}\right)+\dfrac{\sqrt{2}}{2}\cdot\dfrac{2}{3}=\dfrac{-\sqrt{10}+2\sqrt{2}}{6}$.
	\end{itemize}
	} 
\end{ex}

% nguồn: Ok_dungsai_lop11_Lan2
\begin{ex}%[1D1H3-4]%[ma:6550924d]
	Biến đổi biểu thức $\sin a+1$ thành tích. Xét tính đúng sai của các mệnh đề  sau
	\choiceTF 
		{\True $\sin  a+1=2\sin\left(\dfrac{a}{2}+\dfrac{\pi}{4}\right)\cos\left(\dfrac{a}{2}-\dfrac{\pi}{4}\right)$}
		{\True $\sin a+1=\left(\sin\dfrac{a}{2}+\cos\dfrac{a}{2}\right)^2$} 
		{\True $\sin  a+1=2\cos ^2\left(\dfrac{a}{2}-\dfrac{\pi}{4}\right)$}
		{$\sin  a+1=\dfrac{1}{2}{\sin ^2}\left(\dfrac{a}{2}+\dfrac{\pi}{4}\right)$}
	\loigiai{ 
		\begin{itemchoice}
			\itemch \textbf{Đúng.} \\ 
			Vì $\sin a+1=\sin  a+\sin\dfrac{\pi}{2}=2\sin\left(\dfrac{a}{2}+\dfrac{\pi}{4}\right)\cos\left(\dfrac{a}{2}-\dfrac{\pi}{4}\right)$.
			\itemch \textbf{Đúng.} \\ 
			Vì $\sin a+1=2\sin\dfrac{a}{2}\cos\dfrac{a}{2}+\sin ^2\dfrac{a}{2}+\cos^2\dfrac{a}{2}=\left(\sin\dfrac{a}{2}+\cos\dfrac{a}{2}\right)^2$. 
			\itemch \textbf{Đúng.}  \\
			 Vì
			$\begin{array}{l}
			\sin a+1=2\sin\dfrac{a}{2}\cos\dfrac{a}{2}+\sin  ^2\dfrac{a}{2}+\cos^2\dfrac{a}{2}=\left(\sin\dfrac{a}{2}+\cos\dfrac{a}{2}\right)^2=\left[\sqrt2\cos\left(\dfrac{a}{2}-\dfrac{\pi}{4}\right)\right]^2\\
			 =2\cos ^2\left(\dfrac{a}{2}-\dfrac{\pi}{4}\right).
		\end{array}$
		 \itemch \textbf{Sai.} \\
		Vì 
		$\begin{array}{l}
		\sin  a+1=2\sin\dfrac{a}{2}\cos\dfrac{a}{2}+\sin ^2\dfrac{a}{2}+\cos ^2\dfrac{a}{2}=\left(\sin\dfrac{a}{2}+\cos\dfrac{a}{2}\right)^2=\left[\sqrt 2\sin\left(\dfrac{a}{2}+\dfrac{\pi}{4}\right)\right]^2\\
		=2\sin  ^2\left(\dfrac{a}{2}+\dfrac{\pi}{4}\right).
	\end{array}$
	\end{itemchoice}
} 
\end{ex}


%%%%%%%%%%%%%%%% TRẢ LỜI NGẮN %%%%%%%%%%%%%%%%

% nguồn: Ok_TraLoiNgan_Lop11_Lan2
\begin{ex}%[1D1H3-4]%[ma:dbd8338c]
	Rút gọn biểu thức $\dfrac{\sin  7a+\sin 5a}{\sin 7a-\sin 5a}=\tan ma\cot na$. Tính tổng $m+n$.
	\shortans{7} 
	 \loigiai{
		Ta có $\dfrac{\sin 7a+\sin 5a}{\sin 7a-\sin 5a}=\dfrac{2\sin 6a\cos a}{2\cos 6a\sin a}=\tan 6a\cot  a$.\\
		Do đó  $m+n=7$.
	} 
\end{ex}

% nguồn: Ok_TraLoiNgan_Lop11_Lan2
\begin{ex}%[1D1H3-4]%[ma:a8410933]
	Biết $\tan ^2x+\cot  ^2x+2=\dfrac{a}{\sin^22x}$, tìm $a$?
	\shortans{4} 
	\loigiai{ 
		\begin{eqnarray*}
			&&\tan ^2x+\cot ^2x+2=\dfrac{4}{\sin^22x}\\ 
			&\Leftrightarrow&(\tan ^2x+1)+(\cot  ^2x+1)=\dfrac{4}{\sin^22x}\\
			 &\Leftrightarrow&\dfrac{1}{\cos^2x}+\dfrac{1}{\sin^2x}=\dfrac{4}{\sin^22x}\\
			&\Leftrightarrow&\dfrac{\sin^2x+\cos^2x}{\sin^2x{\cos^2}x}=\dfrac{4}{\sin^22x}\\ 
			&\Leftrightarrow&\dfrac{1}{\sin^2x{\cos^2}x}=\dfrac{4}{\sin^22x}\\ 
			 &\Leftrightarrow&\dfrac{1}{\left(\sin x\cos x\right)^2}=\dfrac{4}{\sin^22x}\\
			 &\Leftrightarrow&\dfrac{1}{\dfrac{1}{4}{\sin^2}2x}=\dfrac{4}{\sin^22x}.
		\end{eqnarray*}
		Vậy  $a=4$.
	 }
\end{ex}

% nguồn: Ok_TraLoiNgan_Lop11_Lan2
\begin{ex}%[1D1H3-4]%[ma:5abb7310]
	Biến đổi  thành tổng biểu thức ta được $4\sin 3x\sin 2x\cos x=\cos 2x-\cos 4x-\cos 6x+a$. Khi đó $a$ bằng bao nhiêu?
	 \shortans{1}
	 \loigiai{
		Ta  có
		 \allowdisplaybreaks
		\begin{eqnarray*}
			4\sin 3x\sin 2x\cos x &=&4(\sin  3x\cos x)\sin 2x \\
			&=& 4\cdot\dfrac{1}{2}(\sin  4x+\sin 2x)\sin 2x \\
			&=& 2\sin 4x\sin 2x+2\sin ^22x \\ 
			&=& (\cos 2x-\cos 6x)+2\left(\dfrac{1-\cos 4x}{2}\right)\\ 
			&=&  \cos 2x-\cos 4x-\cos 6x+1.
		\end{eqnarray*}
		 Do đó $a=1$.
	 }
\end{ex}

% nguồn: Ok_TraLoiNgan_Lop11_Lan2
\begin{ex}%[1D1H3-4]%[ma:679dc191]
	Biết rằng $\cos x\cos\left(\dfrac{\pi}{3}-x\right)\cos\left(\dfrac{\pi}{3}+x\right)=k\cos 3x$,  khi đó $k$ bằng bao nhiêu?
	 \shortans{0,25}
	\loigiai{ 
		\allowdisplaybreaks 
		\begin{eqnarray*}
			\cos x\cos\left(\dfrac{\pi}{3}-x\right)\cos\left(\dfrac{\pi}{3}+x\right)  &=& \dfrac{1}{2}\cos x\left(\cos\dfrac{2\pi}{3}+\cos 2x\right)\\
			&=& \dfrac{1}{2}\cos x\cos 2x-\dfrac{1}{4}\cos x=\dfrac{1}{4}(\cos 3x+\cos x)-\dfrac{1}{4}\cos  x\\
			 &=& \dfrac{1}{4}\cos 3x.
		\end{eqnarray*}
		Do đó $k=0{,}25$. 
	 }
\end{ex}

% nguồn: Ok_TraLoiNgan_Lop11_Lan2
\begin{ex}%[1D1H3-4]%[ma:da64e0c6]
	Biết rằng  $\sin 5x-2\sin x(\cos 4x+\cos 2x)=k\sin x$. Khi đó $k$ bằng bao nhiêu?
	\shortans{1} 
	 \loigiai{
		\allowdisplaybreaks 
		\begin{eqnarray*}
			 \sin 5x-2\sin x(\cos 4x+\cos 2x) &=& \sin 5x-2\sin x\cos 4x-2\sin x\cos 2x\\
			&=& \sin 5x-(\sin 5x-\sin 3x)-(\sin 3x-\sin x)\\ 
			&=& \sin  x.
		\end{eqnarray*}
		Do đó $k=1$. 
	} 
\end{ex}

% nguồn: Ok_TraLoiNgan_Lop11_Lan2
\begin{ex}%[1D1H3-4]%[ma:04103700]
	Biết  $\sin ^6x+\cos ^6x=a+\dfrac{3}{8}\cos 4x$, tìm $a$ (kết quả làm tròn đến chữ số thứ 2 sau dấu phẩy)?
	\shortans{0,38} 
	 \loigiai{
		 Ta có
		\begin{eqnarray*}
			\sin ^6x+\cos ^6x&=&\left(\sin^2x+\cos^2x\right)^2-3\sin ^2x{\cos ^2}x\left(\sin^2x+\cos^2x\right)\\ 
			 &=&1-3\sin ^2x{\cos ^2}x\\
			&=&1-\dfrac{3}{4}{\sin ^2}2x\\ 
			 &=&1-\dfrac{3}{8}(1-\cos 4x)\\
			 &=&\dfrac{5}{8}+\dfrac{3}{8}\cos 4x.
		\end{eqnarray*}
	} 
\end{ex}


%%%%%%%%%%%%%%%% TỰ LUẬN %%%%%%%%%%%%%%%%

% nguồn: Ok_TuanLuan_Toan11_Lan2
\begin{ex}%[1D1H3-4]%[ma:0ae08eff]
	 Rút gọn các biểu thức sau:
	\begin{listEX}[2]
		\item $A = \dfrac{\sin \dfrac{\pi}{5} - \sin \dfrac{2\pi}{15}}{\cos \dfrac{\pi}{5} -  \cos \dfrac{2\pi}{15}}$.
		\item $B = \sin \dfrac{\pi}{9} - \sin \dfrac{5\pi}{9} + \sin  \dfrac{7\pi}{9}$.
	\end{listEX}
	 \loigiai{
		\begin{listEX}
			\item $A = \dfrac{\sin \dfrac{\pi}{5} - \sin \dfrac{2\pi}{15}}{\cos \dfrac{\pi}{5} - \cos \dfrac{2\pi}{15}} = \dfrac{2\cos \dfrac{1}{2}\left(\dfrac{\pi}{5} + \dfrac{2\pi}{15}\right) \sin \dfrac{1}{2}\left(\dfrac{\pi}{5} - \dfrac{2\pi}{15}\right)}{-2\sin \dfrac{1}{2}\left(\dfrac{\pi}{5} + \dfrac{2\pi}{15}\right) \sin \dfrac{1}{2}\left(\dfrac{\pi}{5} - \dfrac{2\pi}{15}\right)}= \dfrac{\cos  \dfrac{\pi}{6}}{-\sin \dfrac{\pi}{6}} = -\cot \dfrac{\pi}{6} = -\sqrt{3}$.
			\item $B = \sin \dfrac{\pi}{9} - \sin  \dfrac{5\pi}{9} + \sin \dfrac{7\pi}{9}= 2\sin \dfrac{4\pi}{9} \cos \dfrac{\pi}{3} - \sin \dfrac{5\pi}{9}$
			$= \sin \dfrac{4\pi}{9} - \sin \dfrac{5\pi}{9}$ 
			$= \sin \dfrac{4\pi}{9} - \sin \dfrac{4\pi}{9} = 0$. 
		\end{listEX}
	 }
\end{ex}

% nguồn: Ok_TuanLuan_Toan11_Lan2
\begin{ex}%[1D1H3-4]%[ma:6d196e46]
	Rút gọn biểu  thức sau:
	\begin{listEX}
		\item $A = \dfrac{\cos a + 2\cos 2a + \cos 3a}{\sin a + 2\sin 2a +  \sin 3a}$.
		\item $B = \dfrac{\cos \left(a + \dfrac{\pi}{3}\right) + \cos \left(a - \dfrac{\pi}{3}\right)}{\cot a -  \cot \dfrac{a}{2}}$.
		\item $C = \cos a + \cos (a+b) + \cos (a+2b) +  \ldots + \cos (a+nb) \quad (n \in \mathbb{N})$.
	\end{listEX}
	\loigiai{ 
		\begin{listEX}
			\item $A = \dfrac{(\cos a + \cos 3a) + 2\cos 2a}{(\sin a + \sin 3a) + 2\sin 2a}= \dfrac{2\cos 2a \cos a + 2\cos 2a}{2\sin 2a \cos a + 2\sin 2a} = \dfrac{2\cos 2a (\cos a +  1)}{2\sin 2a (\cos a + 1)} = \cot 2a$.
			\item Ta có $\cos \left(a + \dfrac{\pi}{3}\right) +  \cos \left(a - \dfrac{\pi}{3}\right)= 2\cos a \cos \dfrac{\pi}{3} = \cos a$ và\\
			$\cot a - \cot \dfrac{a}{2} = \dfrac{\cos a}{\sin a} - \dfrac{\cos \dfrac{a}{2}}{\sin \dfrac{a}{2}} = \dfrac{\sin \dfrac{a}{2}\cos a - \cos \dfrac{a}{2} \sin a}{\sin a \sin \dfrac{a}{2}} = \dfrac{\sin \left(\dfrac{a}{2} - a\right)}{\sin a \sin \dfrac{a}{2}} =  \dfrac{-\sin \dfrac{a}{2}}{\sin a \sin \dfrac{a}{2}} = -\dfrac{1}{\sin a}$.\\
			Suy ra $B = \dfrac{\cos a}{-\dfrac{1}{\sin a}} = -\sin a \cos a =  -\dfrac{\sin 2a}{2}$.
			 \item Ta có\\
			 $C \cdot 2\sin \dfrac{b}{2} = 2\sin \dfrac{b}{2} \cos a + 2\sin \dfrac{b}{2}\cos (a+b) + 2\sin \dfrac{b}{2} \cos (a+2b) + \ldots + 2\sin \dfrac{b}{2} \cos (a+nb)$.\\
			 $=\sin\left(\dfrac{b}{2}+a\right)+\sin\left(\dfrac{b}{2}-a\right)+\sin\left(\dfrac{3b}{2}+a\right)+\sin\left(-\dfrac{b}{2}-a\right)+\sin\left(\dfrac{5b}{2}+a\right)+\sin\left(-\dfrac{3b}{2}-a\right)+\ldots+\sin\left(\dfrac{(2n+1)b}{2}+a\right)+\sin\left(-\dfrac{(2n-1)b}{2}-a\right)$.\\
			$=\sin\left(\dfrac{b}{2}-a\right)+\sin\left(\dfrac{(2n+1)b}{2}+a\right)=2\sin \dfrac{(n+1)b}{2} \cdot  \cos\dfrac{nb+2a}{2}$.\\
			Suy ra $C = \dfrac{\sin \dfrac{(n+1)b}{2}  \cdot \cos\dfrac{nb+2a}{2}}{\sin\dfrac{b}{2}}$.
		\end{listEX}
	 }
\end{ex}
